\section{A quantitative near-zero remainder for the HKO cuts}
\label{sec:hko-effective}

The derivative argument can be made effective without enumerating all
twenty-row permutations. Its numerical step is deliberately tiny and
illustrates quantitative control of the true capacity. The separate exact
computation in Theorem~\ref{thm:hko-finite-improvement} gives a substantially
larger ratio increase at a moderate parameter.

\begin{proposition}[Effective improvement]\label{prop:hko-effective}
For the original family~\eqref{eq:ha-family}, with all ten unit-speed cuts,
\[
 \frac{\operatorname{sys}(H_t)}{\operatorname{sys}(H)}
       >1+\frac{t}{200},\qquad 0<t\leq10^{-30}.
\]
\end{proposition}

We give the analytic estimates, including the uniform pivot and the passage
from near-maximizers to exact base maximizers. The finite primitive action
gap and rational constant comparisons have an exact replay; see
\path{certificates/hko-ascent/README.md}.

Put $q(t)=1/(2c_{\EHZ}(H_t))$, $q_0=\tau/5$, and
$V_0=\Vol_4(H)$. The twenty presentation rows have the form
$a_r+tD_r$: one copy of each base normal has $D_r=0$, and its other copy
has the velocity in~\eqref{eq:ha-directions}. A simple word in this bank
therefore contains at most two occurrences of each base normal. For an
ordered word, let
\[
 B_r(t)=(a_r+tD_r,1),\qquad e=(0,0,0,0,1)^{\mathsf T},
 \qquad Q_t(\beta)=\tfrac12\beta^{\mathsf T}G_t\beta,
\]
where $G_t$ is the symmetric ordered-action matrix. Write $D$ for the
matrix whose columns are $(D_r,0)$. Feasibility means
$B(t)\beta=e$ and $\beta\geq0$. The Haim--Kislev formula supplies a
maximizing simple word for every $t$, while $H_t\subset H$ gives
$q(t)\geq q_0$. The all-word response bound proved in
Lemma~\ref{lem:ha-repeat} and Proposition~\ref{prop:ha-certificate} implies
that every exactly maximizing base word with a feasible smooth
continuation has action derivative at most
\begin{equation}\label{eq:he-base-derivative}
 g_*=\frac{2q_0}{125}.
\end{equation}
Indeed, the midpoint identity~\eqref{eq:mc-midpoint} makes that derivative
$2q_0\sum_r\beta_rD_r\cdot x_r^{\rm mid}$, bounded by $2q_0$ times the
total radial-hinge response. Repeated occurrences are covered by the
proved reduction, not by an assumption that all maximizers are simple.
The exact primitive calculation also gives
\begin{equation}\label{eq:he-primitive-gap}
 \Delta=\frac{-\tau+3\tau^3}{20}>\frac1{50}
\end{equation}
as the minimum positive loss from $q_0$ among the $3000$ triangle-pair
cyclic six-orders. These finite predicates are checked in the companion;
the following implications are analytic.

\begin{lemma}[A positive pivot for moving near-maximizers]
\label{lem:he-pivot}
Let $0\leq t\leq10^{-3}$, and let a feasible nonnegative simple word on the
twenty-row bank have $Q_t(\beta)\geq q_0/2$. Five positive coordinates
$I$ have weights greater than $1/4500$, and their augmented base columns
form an invertible matrix $B_I(0)$ satisfying
\[
 |\det B_I(0)|>\tfrac12,\qquad
 \|B_I(0)^{-1}\|_2<66,\qquad
 \|B_I(0)^{-1}\|_\infty<200.
\]
\end{lemma}

\begin{proof}
Let $N=\sqrt5-1=|a_r|$. Both retained and moved rows obey
$|D_r|\leq\sqrt{1+(3N/5)^2}<5/4$, while $N<5/4$.
Bilinearity and $\sum_{r<s}\beta_r\beta_s\leq1/2$ give
\[
 |Q_t(\beta)-Q_0(\beta)|
 \leq\frac{25}{16}t+\frac{25}{32}t^2
 <\frac1{60}<\frac{q_0}{6}.
\]
Here $Q_0(\beta)$ uses the base normals but need not have base closure.
Thus $Q_0(\beta)>q_0/3>1/30$. If $s$ is the total $q$-factor mass,
write $Q_0=s(1-s)B(x,y)$ for normalized occurrence distributions in
the two factors. Each cross pairing has absolute value at most
$N^2<25/16$, so $|B(x,y)|\leq25/16$ and
\[
 s,1-s\geq s(1-s)>\frac8{375}.
\]

Perturbed closure gives the base residual directly, including the
retained rows with zero velocity:
\[
 \left\|\sum_r\beta_ra_r\right\|
 =t\left\|\sum_r\beta_rD_r\right\|
\leq\tfrac54t\leq2t.
\]
After division by either factor mass, the mean residual of its
normalized pentagon-direction distribution is at most
$2t/(8/375)\leq3/32<1/8$. Passing from normals of length $N>1$
to unit normals only decreases it.

Aggregate each normalized factor distribution by its five normal
directions and order those masses decreasingly. Let $r$ be the sum of
the three smallest. The resultant of the two largest unit directions
has norm at least $(1-r)\cos72^\circ>(1-r)/4$; the other three
contribute at most $r$ in the opposite direction. A mean norm below
$1/8$ therefore forces $r>1/16$. Each of the three largest masses is
at least $r/3>1/48$. Multiplying by the factor mass gives aggregate
weight greater than $1/2250$ on each chosen direction. At most two row
occurrences represent one direction, so one occurrence has weight
greater than $1/4500$.

Choose two distinct $q$ directions and three distinct $p$ directions
among these. Their five augmented columns are independent. The
absolute determinant is the product of a two-normal planar determinant
$N^2|\sin(k\pi/5)|>1/2$ and a three-normal affine determinant, which
has one of the exact values $(5\tau+\tau^3)/2$ or $5\tau-\tau^3$,
both greater than one. Thus the five-column determinant exceeds
$1/2$. Each augmented column has norm
$\sqrt{1+N^2}=\sqrt{7-2\sqrt5}<8/5$. Hadamard's inequality bounds
each cofactor by $(8/5)^4<6.56$ and the adjugate's Frobenius norm by
$32.8$. Dividing by the determinant yields
$\|B_I(0)^{-1}\|_2<66$ and hence
$\|B_I(0)^{-1}\|_\infty<200$.
The argument is independent of word order and includes zero-velocity
old rows. For completeness, perturbing these five columns by
$tD_I$ changes their determinant by at most
$5(5/4)t(1.61)^4<1/8$ on the stated interval; the moving minor also
stays invertible.
\end{proof}

\begin{proof}[Proof of Proposition~\ref{prop:hko-effective}]
Fix $0<t\leq10^{-30}$ and a simple word attaining $q(t)$, with weights
$\beta_t$. It satisfies the lemma because $q(t)\geq q_0$.
Set $M=200$, and keep the nonpivot weights $z=(\beta_t)_J$ fixed.
Solve the five pivot weights for $0\leq s\leq t$ by
\begin{equation}\label{eq:he-chart}
 \beta_I(s)=B_I(s)^{-1}\bigl(e-B_J(s)z\bigr),\qquad
 \beta_J(s)=z.
\end{equation}
Let $\beta_0=\beta(0)$. Since $B(0)\beta_t=e-tD\beta_t$,
\[
 \beta_{0,I}-(\beta_t)_I
       =tB_I(0)^{-1}D\beta_t.
\]
The last row of $D$ is zero. Using $|D_r|<2$ gives an infinity-norm
change at most $2Mt$. Hence every pivot of $\beta_0$ exceeds
$1/4500-400\cdot10^{-30}>c:=10^{-4}$. The other weights are unchanged;
$\beta_0$ is a base-feasible probability vector on the same word.

Here are uniform chart bounds, valid more generally for
$0\leq s\leq10^{-8}$. The infinity norm of the five-column
perturbation matrix is at most ten. The Neumann bound gives
$\|B_I(s)^{-1}\|_\infty\leq2M$. The exact difference and differentiated
constraint equations are
\[
 \begin{aligned}
 \beta_I(s)-\beta_{0,I}&=-sB_I(s)^{-1}D\beta_0,\\
 \beta'_I(s)&=-B_I(s)^{-1}D\beta(s),\\
 \beta''_I(s)&=-2B_I(s)^{-1}D_I\beta'_I(s).
 \end{aligned}
\]
Thus $\|\beta(s)\|_1\leq1+20Ms<2$,
$\|\beta'(s)\|_1\leq40M$, and
$\|\beta''(s)\|_1\leq1600M^2$. For the last bound, each
coordinate of $\beta'_I$ has magnitude at most $8M$, each
coordinate of $\beta''_I$ at most $320M^2$, and there are five.
These estimates concern the algebraic chart; its weights need not
stay positive throughout the larger auxiliary interval.

Put $F(s)=Q_s(\beta(s))$. Each moving normal has norm below two there,
so entries of $G_s,G'_s,G''_s$ are bounded in absolute value by
$4,8,8$, respectively. Differentiation gives
\[
\begin{aligned}
 |F''(s)|
 &\leq4\bigl[(40M)^2+2(1600M^2)\bigr]\\
 &\quad+2\cdot8\cdot(40M)\cdot2+\tfrac12\cdot8\cdot2^2\\
 &=19200M^2+1280M+16<2\cdot10^5M^2.
\end{aligned}
\]
With $R=10^5M^2=4\cdot10^9$, Taylor's theorem yields
\begin{equation}\label{eq:he-taylor}
 q(t)=F(t)\leq Q_0(\beta_0)+t\,g(\beta_0)+Rt^2,
\end{equation}
where $g$ is the derivative at zero of the fixed-pivot algebraic
chart. For any base-feasible probability vector $\beta$ on this word,
\[
 g(\beta)=\tfrac12\beta^{\mathsf T}G'_0\beta
                 +(G_0\beta)\cdot u(\beta),\qquad
 u_I(\beta)=-B_I(0)^{-1}D\beta,\quad u_J=0.
\]
The bounds $\|u(\beta)\|_1\leq10M$,
$\max|G_0|\leq4$, and $\max|G'_0|\leq8$ imply
\begin{equation}\label{eq:he-lipschitz}
 |g(\beta)|\leq4+40M\leq44M,\qquad
 |g(\beta)-g(\beta')|\leq88M\|\beta-\beta'\|_1
\end{equation}
for base-feasible probability vectors $\beta,\beta'$ on this word.
The direct quadratic term contributes at most $8\|\beta-\beta'\|_1$;
the other term contributes at most $80M\|\beta-\beta'\|_1$.

Let $d=q_0-Q_0(\beta_0)\geq0$. From $q(t)\geq q_0$,
\eqref{eq:he-taylor}, and~\eqref{eq:he-lipschitz},
\begin{equation}\label{eq:he-small-deficit}
 d\leq44Mt+Rt^2<10^{-20},
 \qquad Q_0(\beta_0)>q_0/2.
\end{equation}
We now round $\beta_0$ to an exact maximizing base word without
changing its order or losing the positive pivot.

Let $s$ be its $q$-factor mass and let $u,v$ be its normalized
factor weights. Each occurrence-wise planar closure polytope has at
most forty vertices: each vertex uses one of five closing triangles
and one of at most two occurrences of each of its three normals.
The same count holds for a subword, which defines a face of the
full occurrence polytope. Write convex decompositions
$u=\sum_i\lambda_i u_i$ and $v=\sum_j\gamma_jv_j$, with at most forty
vertices per factor. If $q_{ij}$ is the action of triangle pair
$(u_i,v_j)$ at balanced factor mass in this fixed order, put
$S=\sum_{ij}\lambda_i\gamma_jq_{ij}$. Then
\begin{equation}\label{eq:he-deficit}
 Q_0(\beta_0)=4s(1-s)S,\qquad
 d=4q_0(s-\tfrac12)^2+4s(1-s)(q_0-S).
\end{equation}
Every $q_{ij}\leq q_0$, and every strict loss is at least
$\Delta$ from~\eqref{eq:he-primitive-gap}. As
$Q_0(\beta_0)>q_0/2$, we have $4s(1-s)>1/2$.
If $\varepsilon$ is the total product mass of nonwinning pairs,
\begin{equation}\label{eq:he-bad-mass}
 \varepsilon\leq\frac{q_0-S}{\Delta}\leq\frac{2d}{\Delta},
 \qquad 2|s-\tfrac12|\leq\sqrt{d/q_0}.
\end{equation}

For $\varepsilon>0$, delete from each marginal every vertex whose
coefficient is less than $2\sqrt\varepsilon$ and renormalize.
At most forty vertices are deleted per marginal, so the discarded
mass is at most $80\sqrt\varepsilon<1$ by
\eqref{eq:he-small-deficit}. A surviving bad pair would have product
mass at least $4\varepsilon$, contradicting the definition of
$\varepsilon$. Thus all surviving pairs win. If $\varepsilon=0$,
keep the original positive supports. Denote the resulting mixtures
by $u^*,v^*$ and put $\beta^*=(u^*/2,v^*/2)$ on the same ordered
occurrences. It is base-feasible and $Q_0(\beta^*)=q_0$.
Each marginal changes by at most $160\sqrt\varepsilon$ in
$\ell^1$, so~\eqref{eq:he-bad-mass} gives
\begin{equation}\label{eq:he-rounding}
 \|\beta_0-\beta^*\|_1
 \leq 2|s-\tfrac12|
       +\tfrac12\|u-u^*\|_1+\tfrac12\|v-v^*\|_1
 \leq\left(\frac1{\sqrt{q_0}}
             +160\sqrt{\frac2\Delta}\right)\sqrt d
 \leq2000\sqrt d.
\end{equation}
By~\eqref{eq:he-small-deficit}, this is less than
$2\cdot10^{-7}<c/2$. Every selected pivot weight of $\beta^*$
is therefore positive. The same fixed-pivot chart gives a feasible
smooth continuation through this exact base maximizer.
Its derivative obeys~\eqref{eq:he-base-derivative}, and
\eqref{eq:he-lipschitz}--\eqref{eq:he-rounding} give
\[
 g(\beta_0)\leq g_*+A\sqrt d,\qquad
 A=88M\cdot2000=35\,200\,000.
\]
This applies the response bound only to an exact maximizing word;
the near-maximizer is handled by the quantified rounding.

Substituting into~\eqref{eq:he-taylor} and maximizing
$-d+At\sqrt d$ over $d\geq0$ by completing the square gives a
bound for the \emph{true global} action:
\begin{equation}\label{eq:he-global-action}
 q(t)\leq q_0+g_*t+(A^2/4+R)t^2
       \leq q_0+\frac{2q_0}{125}t+10^{15}t^2.
\end{equation}
Indeed $A^2/4+R=309\,764\,000\,000\,000<10^{15}$.
No selected feasible branch is substituted for the global optimum.

For volume, on an old facet the gauge of $H_t$ is at least
$1+tf_i$, because the old and moved rows both occur. Other
inequalities only increase the gauge. The normalized cone-volume
measure $\nu$ of~\eqref{eq:mc-nu} and the radial volume formula yield
\[
 \frac{\Vol_4(H_t)}{V_0}
 \leq\int(1+tf)^{-4}\,d\nu
 \leq1-4\nu[f]t+\frac85t^2.
\]
Here $0\leq f_i\leq2/5$, and the second derivative of
$(1+u)^{-4}$ is at most twenty for $u\geq0$.
Put $a=2/125$, $v=4\nu[f]$, and use~\eqref{eq:ha-volume}:
\[
 \eta:=v-2a=\frac{20+64\sqrt5}{9375}>\frac1{100}.
\]
As $q_0>1/10$, \eqref{eq:he-global-action} gives
\[
 \left(\frac{q(t)}{q_0}\right)^2
 \leq(1+at+10^{16}t^2)^2
 \leq1+2at+3\cdot10^{16}t^2.
\]
Multiplying this nonnegative bound by the volume bound and
discarding only negative cross terms gives
\[
 \frac{q(t)^2\Vol_4(H_t)}{q_0^2V_0}
 \leq1-\eta t+4\cdot10^{16}t^2
 <1-\frac{t}{200},
\]
because $4\cdot10^{16}t<1/200$ on the stated interval.
Finally $\operatorname{sys}(H_t)=1/(8q(t)^2\Vol_4(H_t))$,
so taking reciprocals proves the proposition.
\end{proof}

At $t=10^{-30}$ this gives only the certified lower bound
$5\cdot10^{-33}$ on relative gain. The reusable ingredient is the
uniform quadratic remainder for the global optimal action; determining
the maximum along the family remains a separate problem.
